% ===================== 第二阶段：技巧进阶与伴奏应用（第 11-20 课） =====================
\pdfbookmark[1]{第二阶段 · 技巧进阶与伴奏应用}{stg2}
\stagepage{2}{stageB}{技巧进阶与伴奏应用}{第 11--20 课}%
{掌握击弦、勾弦、滑弦、闷音、强力和弦、推弦与揉弦等核心技巧，完成十二个自然大调的调内和弦推导，并以《One》完成伴奏阶段的结课曲目练习。}%
{\stgitem{第 11 课 · 击弦、勾弦、滑弦与和弦内外音}
 \stgitem{第 12 课 · 击勾、双滑与断音练习}
 \stgitem{第 13 课 · 提速练习与和弦外音}
 \stgitem{第 14 课 · 调式学习：十二个自然大调的和弦推导}
 \stgitem{第 15 课 · 和弦外音与小拇指击勾弦}}%
{\stgitem{第 16 课 · 闷音、强力和弦与滑奏}
 \stgitem{第 17 课 · 阶段练习课（二）}
 \stgitem{第 18 课 · 阶段练习课（三）}
 \stgitem{第 19 课 · 制音扫弦、切音、推弦与揉弦}
 \stgitem{第 20 课 · 伴奏阶段结课练习}}
\SetStage{stageB}{第二阶段 · 技巧进阶与伴奏应用}{第 11--20 课}

\pdfbookmark[2]{第 11 课 · 击弦、勾弦、滑弦与和弦内外音}{lsn11}
\begin{lessoncard}{11}{击弦、勾弦、滑弦与和弦内外音}
\cardsec{\faBullseye}{教学目标}
\begin{itemize}
\item 掌握击弦、勾弦和滑弦三种左手连奏技巧的基本动作。
\item 能够在保持和弦手型的情况下，完成和弦内音之间的击勾弦装饰。
\item 能够把击弦、勾弦与滑弦连接在同一个短句中。
\item 初步区分和弦内音与和弦外音的听感及使用方式。
\end{itemize}
\cardsec{\faIcon{book-open}}{教学内容}
\begin{itemize}
\item 击弦、勾弦和滑弦的发力位置、动作幅度与发音要求。
\item 在和弦手型中保留主要按弦手指，并加入和弦内的击勾弦变化。
\item 将击弦、勾弦和滑弦按不同顺序组合，形成连贯的装饰性短句。
\item 对比和弦内音与和弦外音，理解外音需要解决或回到稳定音的基本听感。
\item 练习过程中保持左手放松，避免只靠大幅度甩动换取音量。
\end{itemize}
\cardsec{\faIcon{pen-fancy}}{课堂练习}
\begin{itemize}
\item 在单弦上分别完成击弦、勾弦和滑弦的基础练习。
\item 保持一个基础和弦手型，在和弦音之间加入击弦和勾弦。
\item 设计一至两小节短句，把击弦、勾弦和滑弦三种技巧融合起来循环练习。
\item 对比演奏和弦内音与和弦外音，听辨稳定、紧张与解决的差别。
\item 检查每个连奏音的音量是否均衡，是否出现杂音、断音或动作过大。
\end{itemize}
\begin{fignote}
\begin{center}
% 击弦（Hammer-on）示意
\begin{tikzpicture}[x=0.9cm,y=0.5cm]
  \foreach \s in {0,1,2}{\draw[inkGray!45, line width=0.5pt] (0,\s) -- (4,\s);}
  \foreach \f in {0,...,4}{\draw[inkGray, line width=0.9pt] (\f,-0.35) -- (\f,2.35);}
  \fill[stageB] (0.5,1) circle (0.15);
  \fill[stageB!55] (2.5,1) circle (0.15);
  \draw[stageB, line width=1.1pt, -{Stealth[length=2.2mm]}]
    (0.65,1.5) .. controls (1.3,2.5) and (2.0,2.5) .. (2.5,1.35);
  \node[font=\footnotesize\sffamily, color=inkGray, anchor=north] at (2,-0.9) {击弦：手指敲落到高品};
\end{tikzpicture}\hspace{12mm}%
% 勾弦（Pull-off）示意
\begin{tikzpicture}[x=0.9cm,y=0.5cm]
  \foreach \s in {0,1,2}{\draw[inkGray!45, line width=0.5pt] (0,\s) -- (4,\s);}
  \foreach \f in {0,...,4}{\draw[inkGray, line width=0.9pt] (\f,-0.35) -- (\f,2.35);}
  \fill[stageB] (2.5,1) circle (0.15);
  \fill[stageB!55] (0.5,1) circle (0.15);
  \draw[stageB, line width=1.1pt, -{Stealth[length=2.2mm]}]
    (2.35,1.5) .. controls (1.7,2.5) and (1.0,2.5) .. (0.55,1.35);
  \node[font=\footnotesize\sffamily, color=inkGray, anchor=north] at (2,-0.9) {勾弦：手指勾离回到低品};
\end{tikzpicture}\hspace{12mm}%
% 滑弦（Slide）示意
\begin{tikzpicture}[x=0.9cm,y=0.5cm]
  \foreach \s in {0,1,2}{\draw[inkGray!45, line width=0.5pt] (0,\s) -- (4,\s);}
  \foreach \f in {0,...,4}{\draw[inkGray, line width=0.9pt] (\f,-0.35) -- (\f,2.35);}
  \fill[stageB] (0.5,1) circle (0.15);
  \fill[stageB!55] (3.5,1) circle (0.15);
  \draw[stageB, line width=1.6pt, -{Stealth[length=2.2mm]}] (0.75,1) -- (3.2,1);
  \node[font=\footnotesize\sffamily, color=inkGray, anchor=north] at (2,-0.9) {滑弦：按弦沿弦滑向目标品};
\end{tikzpicture}
\figcaption{三种左手连奏技巧示意（深色为起始音，浅色为目标音）}
\end{center}
\begin{fignotes}
\item 击弦：手指像小锤一样敲落到高品位发音，只拨第一个音，第二个音靠敲击发声。
\item 勾弦：手指向掌心方向勾离琴弦回到低品位，勾的瞬间带出下一个音。
\item 滑弦：按实后手指不离弦、保持压力沿弦滑到目标品；三种技巧都要求音量均衡、无杂音。
\end{fignotes}
\end{fignote}
\end{lessoncard}

\pdfbookmark[2]{第 12 课 · 击勾、双滑与断音练习}{lsn12}
\begin{lessoncard}{12}{击勾、双滑与断音练习}
\songrow{选曲}{《Fade to Black》}
\cardsec{\faBullseye}{教学目标}
\begin{itemize}
\item 巩固击弦、勾弦和滑弦的发音稳定性。
\item 掌握击勾、双滑与断音之间的组合与衔接。
\item 能够在曲目片段中识别并完成相应技巧。
\item 提高左、右手在连奏与断音切换时的协调性。
\end{itemize}
\cardsec{\faIcon{book-open}}{教学内容}
\begin{itemize}
\item 击弦与勾弦的连续组合，以及连奏音之间的力度均衡。
\item 上行、下行及往返双滑的动作与目标音控制。
\item 断音的基本处理，包括左手放松止音与右手辅助控制。
\item 《Fade to Black》选段中的技巧拆分、分句与连接。
\end{itemize}
\cardsec{\faIcon{pen-fancy}}{课堂练习}
\begin{itemize}
\item 分别循环击勾组合、双滑组合与断音练习。
\item 将击勾、双滑和断音放入一至两小节的综合短句。
\item 拆分练习《Fade to Black》选段中的难点动作。
\item 先保证每个音清晰，再逐步连接完整乐句。
\end{itemize}
\begin{fignote}
\begin{center}
% 双滑示意
\begin{tikzpicture}[x=0.9cm,y=0.5cm]
  \foreach \s in {0,1,2}{\draw[inkGray!45, line width=0.5pt] (0,\s) -- (4,\s);}
  \foreach \f in {0,...,4}{\draw[inkGray, line width=0.9pt] (\f,-0.35) -- (\f,2.35);}
  \fill[stageB] (0.5,1) circle (0.15);
  \fill[stageB!55] (2.5,1) circle (0.15);
  \draw[stageB, line width=1.4pt,
    {Stealth[length=1.9mm]}-{Stealth[length=1.9mm]}] (0.78,1) -- (2.22,1);
  \node[font=\footnotesize\sffamily, color=inkGray, anchor=north] at (2,-0.9) {双滑：滑到目标品再滑回};
\end{tikzpicture}\hspace{14mm}%
% 断音示意
\begin{tikzpicture}[x=0.9cm,y=0.5cm]
  \foreach \s in {0,1,2}{\draw[inkGray!45, line width=0.5pt] (0,\s) -- (4,\s);}
  \foreach \f in {0,...,4}{\draw[inkGray, line width=0.9pt] (\f,-0.35) -- (\f,2.35);}
  \fill[stageB] (1.0,1) circle (0.15);
  \draw[stageB!70, line width=0.9pt] (1.25,1.28) .. controls (1.55,1.5) and (1.8,1.06) .. (2.05,1.28);
  \draw[stageB, line width=1.6pt] (2.45,0.62) -- (2.45,1.38);
  \draw[stageB, line width=1.6pt] (2.62,0.62) -- (2.62,1.38);
  \node[font=\footnotesize\sffamily, color=inkGray, anchor=north] at (2,-0.9) {断音：发声后立即止住};
\end{tikzpicture}
\end{center}
\begin{fignotes}
\item 双滑是有去有回的滑弦：滑到目标品站稳发音，再按原路滑回，两端音都要清楚。
\item 断音靠「止」不靠「砍」：左手轻抬放松（不离弦）或右手掌侧搭弦，把余音干净掐断。
\item 连奏与断音交替时最考验双手配合：先分开循环，再放进同一条短句里衔接。
\end{fignotes}
\end{fignote}
\end{lessoncard}

\pdfbookmark[2]{第 13 课 · 提速练习与和弦外音}{lsn13}
\begin{lessoncard}{13}{提速练习与和弦外音}
\songrow{曲目练习}{谢天笑《循环的太阳》}
\cardsec{\faBullseye}{教学目标}
\begin{itemize}
\item 建立正确的提速练习方法，提高演奏的稳定性与连贯度。
\item 进一步理解和弦外音在旋律和乐句中的作用。
\item 将前面学习的连奏技巧应用到具体曲目中。
\item 在速度增加时保持动作放松、音符清晰和节奏稳定。
\end{itemize}
\cardsec{\faIcon{book-open}}{教学内容}
\begin{itemize}
\item 分段、缩小循环范围和逐步提速的练习方法。
\item 提速过程中左手动作幅度、右手拨弦与双手同步的检查。
\item 和弦外音的经过、紧张与解决，以及回到和弦内音的听感。
\item 《循环的太阳》相关片段的技术拆解与连接。
\end{itemize}
\cardsec{\faIcon{pen-fancy}}{课堂练习}
\begin{itemize}
\item 选择短小乐句进行慢速准确练习，再逐级提高速度。
\item 单独循环容易卡顿或双手不同步的位置。
\item 在固定和弦背景下加入简单和弦外音，并回到稳定音。
\item 分句练习《循环的太阳》，再逐步连接成完整段落。
\end{itemize}
\begin{fignote}
\begin{center}
\begin{tikzpicture}[line cap=round]
  \draw[inkGray, line width=1.1pt, fill=stagecur!8]
    (-1.15,0) -- (-0.45,3.2) -- (0.45,3.2) -- (1.15,0) -- cycle;
  \draw[inkGray, line width=1.1pt, fill=stagecur!15, rounded corners=1pt]
    (-1.28,-0.34) rectangle (1.28,0);
  \draw[inkGray!65, line width=0.8pt] (0,0.55) -- (0,2.85);
  \foreach \y in {0.85,1.2,1.55,1.9,2.25,2.6}{
    \draw[inkGray!65, line width=0.7pt] (-0.13,\y) -- (0.13,\y);}
  \draw[stagecur, line width=1.5pt, rotate around={17:(0,0.5)}] (0,0.5) -- (0,2.95);
  \begin{scope}[rotate around={17:(0,0.5)}]
    \fill[stagecur, rounded corners=0.8pt] (-0.17,1.95) rectangle (0.17,2.32);
  \end{scope}
  \draw[stagecur!55, line width=0.8pt, dashed,
    {Stealth[length=1.6mm]}-{Stealth[length=1.6mm]}]
    (-0.95,2.98) arc[start angle=112, end angle=68, radius=2.55];
\end{tikzpicture}
\end{center}
\begin{fignotes}
\item 提速前先分段、缩小循环范围：把最容易卡顿的一两小节单独拎出来循环。
\item 从能完全弹对的慢速开始，每次只加一小格速度，弹稳后再进下一档；一卡顿就退回上一档。
\item 加速过程中随时检查：左手动作幅度有没有变大、双手是否依然同步、身体是否放松。
\end{fignotes}
\end{fignote}
\end{lessoncard}

\pdfbookmark[2]{第 14 课 · 调式学习：十二个自然大调的和弦推导}{lsn14}
\begin{lessoncard}{14}{调式学习：十二个自然大调的和弦推导}
\cardsec{\faBullseye}{教学目标}
\begin{itemize}
\item 理解自然大调音阶与调内和弦之间的关系。
\item 掌握自然大调各级三和弦的固定性质。
\item 能够从任意大调音阶推导出该调的全部调内三和弦。
\item 建立十二个自然大调之间可以使用同一套推导逻辑的认识。
\end{itemize}
\cardsec{\faIcon{book-open}}{教学内容}
\begin{itemize}
\item 自然大调音阶的音程排列及七个音级。
\item 以音阶各级音为根音，按三度叠置构成三和弦的方法。
\item 大调体系中各级和弦性质：I、IV、V 级为大三和弦，ii、iii、vi 级为小三和弦，vii\textdegree{} 级为减三和弦。
\item 使用同一推导方法完成十二个自然大调的调内和弦整理。
\item 调内和弦级数与实际和弦名称之间的对应。
\end{itemize}
\cardsec{\faIcon{pen-fancy}}{课堂练习}
\begin{itemize}
\item 先以一个熟悉的大调完整示范音阶与和弦推导。
\item 给出不同主音，由学生写出自然大调音阶。
\item 根据音阶逐级叠置三度，推导各级三和弦及其性质。
\item 抽取多个调进行口头或书面推导，检查是否真正掌握规律。
\end{itemize}
\begin{fignote}
\begin{center}
\begin{tikzpicture}
  \foreach \x/\n in {1/I, 4/IV, 5/V}{
    \node[fill=stagecur, text=white, rounded corners=1.2mm, minimum width=1.15cm,
      minimum height=0.62cm, font=\small\sffamily\bfseries] at (\x*1.45,0) {\n};}
  \foreach \x/\n in {2/ii, 3/iii, 6/vi}{
    \node[fill=stagecur!25, text=stagecur!60!black, rounded corners=1.2mm,
      minimum width=1.15cm, minimum height=0.62cm, font=\small\sffamily\bfseries]
      at (\x*1.45,0) {\n};}
  \node[draw=stagecur!70, dashed, text=stagecur!70!black, rounded corners=1.2mm,
    minimum width=1.15cm, minimum height=0.62cm, font=\small\sffamily\bfseries]
    at (7*1.45,0) {vii\textdegree};
  \foreach \x/\n in {1/C, 2/Dm, 3/Em, 4/F, 5/G, 6/Am, 7/Bdim}{
    \node[font=\scriptsize\sffamily, inkGray] at (\x*1.45,-0.72) {\n};}
  \node[font=\scriptsize\sffamily, inkGray!80] at (5.8,-1.45)
    {\textcolor{stagecur}{\rule[-1pt]{6pt}{6pt}}\ 大三和弦\quad
     \textcolor{stagecur!25}{\rule[-1pt]{6pt}{6pt}}\ 小三和弦\quad
     \tikz\draw[stagecur!70, dashed, line width=0.6pt] (0,0) rectangle (0.21,0.21);\ 减三和弦};
\end{tikzpicture}
\end{center}
\begin{fignotes}
\item 各级性质固定不变：I、IV、V 级大三和弦，ii、iii、vi 级小三和弦，vii\textdegree{} 级减三和弦。
\item 下排是 C 大调的示范结果：C、Dm、Em、F、G、Am、Bdim——照着级数表逐级叠三度即可推出。
\item 换任何一个大调，级数与性质的对应关系完全相同，所以十二个调共用同一套推导逻辑。
\end{fignotes}
\end{fignote}
\end{lessoncard}

\pdfbookmark[2]{第 15 课 · 和弦外音与小拇指击勾弦}{lsn15}
\begin{lessoncard}{15}{和弦外音与小拇指击勾弦}
\songrow{曲目练习}{《天鹅之死》前奏}
\cardsec{\faBullseye}{教学目标}
\begin{itemize}
\item 加强小拇指的独立性、力度和控制能力。
\item 能够使用小拇指完成清晰的击弦与勾弦。
\item 继续理解和弦外音在乐句中的装饰、经过与解决。
\item 将小拇指击勾弦和和弦外音应用到曲目前奏中。
\end{itemize}
\cardsec{\faIcon{book-open}}{教学内容}
\begin{itemize}
\item 小拇指击弦、勾弦的落点、发力与动作幅度。
\item 小拇指与无名指之间的独立性和配合。
\item 和弦内音、和弦外音在同一短句中的连接。
\item 《天鹅之死》前奏的指法、连奏技巧与分句练习。
\end{itemize}
\cardsec{\faIcon{pen-fancy}}{课堂练习}
\begin{itemize}
\item 进行小拇指专项击弦、勾弦和组合练习。
\item 在固定和弦或和弦音框架内加入经过性的和弦外音。
\item 拆分《天鹅之死》前奏中的小拇指与连奏难点。
\item 对音量偏小、杂音或动作僵硬的位置进行短循环纠正。
\end{itemize}
\begin{fignote}
\begin{center}
\begin{tikzpicture}[x=0.9cm,y=0.5cm]
  \foreach \s in {0,1,2}{\draw[inkGray!45, line width=0.5pt] (0,\s) -- (4,\s);}
  \foreach \f in {0,...,4}{\draw[inkGray, line width=0.9pt] (\f,-0.35) -- (\f,2.35);}
  \fill[stageB] (0.5,1) circle (0.17);
  \node[white,font=\tiny\sffamily] at (0.5,1) {1};
  \fill[stageB!55] (2.5,1) circle (0.17);
  \node[white,font=\tiny\sffamily] at (2.5,1) {4};
  \draw[stageB, line width=1.1pt, -{Stealth[length=2.2mm]}]
    (0.7,1.55) .. controls (1.3,2.55) and (2.0,2.55) .. (2.5,1.4);
  \node[font=\footnotesize\sffamily, color=inkGray, anchor=north] at (2,-0.9)
    {小指击弦：食指按稳，小指敲落};
\end{tikzpicture}
\end{center}
\begin{fignotes}
\item 食指（1）先把低品位按稳当支点，小指（4）像小锤一样垂直敲落到高品位。
\item 小指动作幅度要小：从贴近琴弦的高度直接落下，音量不够就练落点准度，不要抡大臂。
\item 勾弦方向朝掌心：小指勾离时轻带琴弦发音，与无名指分开单独练，防止两指粘连。
\end{fignotes}
\end{fignote}
\end{lessoncard}

\pdfbookmark[2]{第 16 课 · 闷音、强力和弦与滑奏}{lsn16}
\begin{lessoncard}{16}{闷音、强力和弦与滑奏}
\cardsec{\faBullseye}{教学目标}
\begin{itemize}
\item 掌握基础闷音的发音与力度控制。
\item 理解强力和弦的基本结构和常用指型。
\item 能够在不同把位之间完成强力和弦滑奏。
\item 提高重型节奏中的整齐度、颗粒感与连贯性。
\end{itemize}
\cardsec{\faIcon{book-open}}{教学内容}
\begin{itemize}
\item 右手掌根闷音的位置、压力与音色变化。
\item 强力和弦的根音、五度音及常见两音、三音指型。
\item 强力和弦在不同弦组和把位上的移动方法。
\item 滑奏过程中手型保持、目标品定位与杂音控制。
\item 闷音、开放音和滑奏之间的动态与音色对比。
\end{itemize}
\cardsec{\faIcon{pen-fancy}}{课堂练习}
\begin{itemize}
\item 在低音弦上对比开放拨弦、轻闷音与重闷音。
\item 使用固定强力和弦完成重复节奏练习。
\item 在两个或多个把位之间练习强力和弦滑奏。
\item 将闷音、强力和弦与滑奏组合成短小节奏段落。
\end{itemize}
\begin{fignote}
\begin{center}
\begin{chordbox}{E5（两音）}
  \cbopen{6}\cbdot{5}{2}\cbmute{4}\cbmute{3}\cbmute{2}\cbmute{1}
\end{chordbox}\hspace{10mm}%
\begin{chordbox}{A5（三音）}
  \cbmute{6}\cbopen{5}\cbdot{4}{2}\cbdot{3}{2}\cbmute{2}\cbmute{1}
\end{chordbox}\hspace{10mm}%
\photo[3.9cm]{ai/palm_mute.png}
\end{center}
\begin{fignotes}
\item 强力和弦只有根音和五度音，没有三音，所以不分大小、失真音色下依然干净有力。
\item 两音型＝根音＋五度；三音型再叠一个高八度根音，声音更厚；同一指型可整体平移换把。
\item 掌根闷音（右图）：右手掌外侧贴在琴桥正上方的弦上，越靠琴桥越亮、越往前越闷——压力大小决定颗粒感。
\end{fignotes}
\end{fignote}
\end{lessoncard}

\pdfbookmark[2]{第 17 课 · 阶段练习课（二）}{lsn17}
\begin{lessoncard}{17}{阶段练习课（二）}
\songrow{选曲}{《Fade to Black》副歌，或《Du Hast》前奏}
\cardsec{\faBullseye}{教学目标}
\begin{itemize}
\item 综合复习第二阶段已经学习的连奏、闷音、强力和弦及和弦切换技巧。
\item 根据所选曲目集中解决节奏、手型和段落衔接问题。
\item 提高连续演奏时的稳定性与完整度。
\item 建立针对具体问题进行拆分练习的能力。
\end{itemize}
\cardsec{\faIcon{book-open}}{教学内容}
\begin{itemize}
\item 《Fade to Black》副歌中的和弦、节奏与段落连接。
\item 或《Du Hast》前奏中的强力和弦、闷音与节奏控制。
\item 对前面课程中暴露出的技术问题进行集中复习。
\item 本课以曲目应用和纠错为主，不集中加入大量新技巧。
\end{itemize}
\cardsec{\faIcon{pen-fancy}}{课堂练习}
\begin{itemize}
\item 根据学生情况选择《Fade to Black》副歌或《Du Hast》前奏。
\item 分别练习左手和弦或指型、右手节奏，再进行合并。
\item 对容易中断、抢拍或杂音较多的位置进行小节循环。
\item 最后尝试连续完成所选段落，并记录仍需解决的问题。
\end{itemize}
\begin{fignote}
\begin{center}
\begin{adjustbox}{max width=\textwidth}
\stepchip{选定段落}\steparrow
\stepchip{左右手分开练}\steparrow
\stepchip{合并}\steparrow
\stepchip{问题小节循环}\steparrow
\stepchip{连奏＋记录问题}
\end{adjustbox}
\end{center}
\begin{fignotes}
\item 曲目课的固定路径：先拆手（左手指型、右手节奏各自练顺）再合并，合不上就回去重拆。
\item 循环的单位是「小节」不是「整段」：哪一小节断、抢拍、杂音多，就只圈那一小节。
\item 下课前必须连奏一遍并写下没解决的问题——它们就是下节课的起点。
\end{fignotes}
\end{fignote}
\end{lessoncard}

\pdfbookmark[2]{第 18 课 · 阶段练习课（三）}{lsn18}
\begin{lessoncard}{18}{阶段练习课（三）}
\songrow{曲目练习}{谢天笑《把夜晚染黑》前奏分解与 Solo}
\cardsec{\faBullseye}{教学目标}
\begin{itemize}
\item 巩固分解和弦、旋律演奏与 Solo 乐句之间的衔接。
\item 提高右手换弦、左手连贯移动和乐句完整性。
\item 能够分段完成前奏分解与 Solo，并逐步连接。
\item 继续强化曲目中的节奏、音色与杂音控制。
\end{itemize}
\cardsec{\faIcon{book-open}}{教学内容}
\begin{itemize}
\item 《把夜晚染黑》前奏分解部分的右手顺序与左手指型。
\item 前奏分解与 Solo 之间的转换和准备动作。
\item Solo 乐句中的指法、连奏和换把处理。
\item 段落拆分、难点循环与整体连接方法。
\end{itemize}
\cardsec{\faIcon{pen-fancy}}{课堂练习}
\begin{itemize}
\item 单独练习前奏分解的右手顺序和音量均衡。
\item 分句练习 Solo，对换把和连奏位置进行专项处理。
\item 单独循环前奏进入 Solo 的连接位置。
\item 在保证清晰度的前提下完成完整段落演奏。
\end{itemize}
\begin{fignote}
\begin{center}
\begin{adjustbox}{max width=\textwidth}
\stepchip{前奏分解单练}\steparrow
\stepchip{Solo 分句单练}\steparrow
\stepchip{连接位置专项循环}\steparrow
\stepchip{完整段落}
\end{adjustbox}
\end{center}
\begin{fignotes}
\item 前奏和 Solo 是两种手感（分解拨弦 vs 旋律连奏），先各自练熟再谈衔接。
\item 断点几乎都发生在「进入 Solo 的那一下」：把连接前后各一小节单独抠成一个循环。
\item 连接时提前半拍准备左手换把位置，眼睛和手都先到，音再到。
\end{fignotes}
\end{fignote}
\end{lessoncard}

\pdfbookmark[2]{第 19 课 · 制音扫弦、切音、推弦与揉弦}{lsn19}
\begin{lessoncard}{19}{制音扫弦、切音、推弦与揉弦}
\cardsec{\faBullseye}{教学目标}
\begin{itemize}
\item 掌握制音扫弦和切音的基本动作与节奏效果。
\item 能够完成基础推弦，并建立目标音高意识。
\item 掌握基础揉弦动作，使长音具有稳定的音高变化。
\item 区分节奏型技巧与旋律表现技巧的不同作用。
\end{itemize}
\cardsec{\faIcon{book-open}}{教学内容}
\begin{itemize}
\item 左、右手制音在扫弦中的配合，以及杂音控制。
\item 切音的时值控制、放松动作与律动效果。
\item 推弦的发力方向、手指支撑与目标音高。
\item 揉弦的动作方向、幅度、速度与稳定性。
\item 将制音、切音用于节奏，将推弦、揉弦用于旋律乐句。
\end{itemize}
\cardsec{\faIcon{pen-fancy}}{课堂练习}
\begin{itemize}
\item 在单个和弦上交替练习正常扫弦、制音扫弦和切音。
\item 使用短节奏型练习切音位置与停顿长度。
\item 先弹目标音，再进行半音或全音推弦并对照音高。
\item 在长音上练习均匀揉弦，再加入简单旋律短句。
\end{itemize}
\begin{fignote}
\begin{center}
% 推弦示意
\begin{tikzpicture}[x=0.9cm,y=0.5cm,baseline=(current bounding box.center)]
  \draw[inkGray!45, line width=0.5pt] (0,0) -- (4,0);
  \draw[inkGray!45, line width=0.5pt] (0,2) -- (4,2);
  \draw[inkGray!30, line width=0.5pt, dashed] (0,1) -- (4,1);
  \draw[inkGray!70, line width=0.6pt] (0,1) -- (1.1,1)
    .. controls (1.3,1) and (1.4,1.62) .. (1.5,1.65)
    .. controls (1.6,1.62) and (1.7,1) .. (1.9,1) -- (4,1);
  \foreach \f in {0,...,4}{\draw[inkGray, line width=0.9pt] (\f,-0.35) -- (\f,2.35);}
  \fill[stageB] (1.5,1.65) circle (0.16);
  \draw[stageB, line width=1.1pt, -{Stealth[length=2mm]}] (1.5,0.35) -- (1.5,1.32);
  \node[font=\footnotesize\sffamily, color=inkGray, anchor=north] at (2,-0.9) {推弦：顶弦升高音高};
\end{tikzpicture}\hspace{14mm}%
% 揉弦示意
\begin{tikzpicture}[x=0.9cm,y=0.5cm,baseline=(current bounding box.center)]
  \foreach \s in {0,1,2}{\draw[inkGray!45, line width=0.5pt] (0,\s) -- (4,\s);}
  \foreach \f in {0,...,4}{\draw[inkGray, line width=0.9pt] (\f,-0.35) -- (\f,2.35);}
  \fill[stageB] (1.5,1) circle (0.16);
  \draw[stageB, line width=1.1pt, {Stealth[length=1.8mm]}-{Stealth[length=1.8mm]}]
    (1.5,0.25) -- (1.5,1.78);
  \draw[stageB!70, line width=0.8pt]
    (2.1,1.35) .. controls (2.35,1.6) and (2.55,1.1) .. (2.8,1.35);
  \draw[stageB!45, line width=0.8pt]
    (2.1,1.9) .. controls (2.35,2.15) and (2.55,1.65) .. (2.8,1.9);
  \node[font=\footnotesize\sffamily, color=inkGray, anchor=north] at (2,-0.9) {揉弦：按弦后往返摆动};
\end{tikzpicture}
\end{center}
\begin{fignotes}
\item 推弦：按稳后向上顶弦，虚线是弦的原位置；先弹出目标音、再推弦对照音高是否到位。
\item 推弦靠手腕转动发力，食指、中指在推弦指后方一起支撑，不要只用一根手指硬顶。
\item 揉弦：按稳后小幅往返摆动，幅度和速度保持均匀，让长音产生稳定的音高波动。
\end{fignotes}
\end{fignote}
\end{lessoncard}

\pdfbookmark[2]{第 20 课 · 伴奏阶段结课练习}{lsn20}
\begin{lessoncard}{20}{伴奏阶段结课练习}
\songrow{结课曲目}{《One》}
\cardsec{\faBullseye}{教学目标}
\begin{itemize}
\item 综合运用第一、第二阶段学习的和弦、分解、扫弦及技巧内容。
\item 完成伴奏阶段的总结与结课曲目练习。
\item 提高段落记忆、节奏稳定、音色控制和连续演奏能力。
\item 通过完整曲目检查学生仍需巩固的技术环节。
\end{itemize}
\cardsec{\faIcon{book-open}}{教学内容}
\begin{itemize}
\item 《One》所选伴奏段落的结构、指法与右手演奏方式。
\item 不同段落之间的音色、力度和演奏方式变化。
\item 曲目中涉及的分解、和弦、节奏及换把内容。
\item 对第一、第二阶段重点技术进行综合复查。
\end{itemize}
\cardsec{\faIcon{pen-fancy}}{课堂练习}
\begin{itemize}
\item 按段落拆分练习《One》，分别处理左手和右手问题。
\item 单独循环段落连接与容易中断的位置。
\item 尝试连续演奏所学部分，检查节奏、杂音和动作紧张。
\item 完成伴奏阶段总结，明确后续需要长期保留的练习内容。
\end{itemize}
\begin{fignote}
\begin{center}
\photo[4.6cm]{ai/stage_solo.png}
\end{center}
\begin{fignotes}
\item 结课检查三件事：段落记得住、节奏站得稳、音色控得住——对应本课的三个练习重点。
\item 连续演奏时别为一个错音停下来：记住位置继续走，弹完再回头修，这是完整度训练的核心。
\item 结课不是终点：把总结出的薄弱环节写成长期保留的练习清单，带进第三阶段。
\end{fignotes}
\end{fignote}
\end{lessoncard}
